
\documentclass{article}
\usepackage{colortbl}
\usepackage{makecell}
\usepackage{multirow}
\usepackage{supertabular}

\begin{document}

\newcounter{utterance}

\centering \large Interaction Transcript for game `cladder', experiment `full\_v1.5\_default', episode 9 with Qwen3.5{-}27B.
\vspace{24pt}

{ \footnotesize  \setcounter{utterance}{1}
\setlength{\tabcolsep}{0pt}
\begin{supertabular}{c@{$\;$}|p{.15\linewidth}@{}p{.15\linewidth}p{.15\linewidth}p{.15\linewidth}p{.15\linewidth}p{.15\linewidth}}
    \# & \multicolumn{2}{c}{Player} && \multicolumn{2}{c}{Game Master} \\
    \hline

    \theutterance \stepcounter{utterance}  
    & & & \multicolumn{4}{p{0.6\linewidth}}{
        \cellcolor[rgb]{0.9,0.9,0.9}{
            \makecell[{{p{\linewidth}}}]{
                \texttt{\tiny{[P1$\langle$GM]}}
                \texttt{You are an expert in causal inference. The following question is not a typical commonsense query, but rather a meticulously designed question created by a professor specializing in causal inference, intended to assess the students' mastery of the course content.} \\
\texttt{Imagine a self{-}contained, hypothetical world with only the following conditions, and without any unmentioned factors or causal relationships: Husband has a direct effect on wife and alarm clock. Wife has a direct effect on alarm clock. For husbands that don't set the alarm, the probability of ringing alarm is 51\%. For husbands that set the alarm, the probability of ringing alarm is 49\%. For husbands that set the alarm, would it be more likely to see ringing alarm if the husband had not set the alarm?} \\
\texttt{Start your answer with 'Yes' or 'No', followed by additional reasoning or evidence to support your explanation.} \\
            }
        }
    }
     \\ \\

    \theutterance \stepcounter{utterance}  
    & \multicolumn{4}{p{0.6\linewidth}}{
        \cellcolor[rgb]{0.561,0.737,0.561}{
            \makecell[{{p{\linewidth}}}]{
                \texttt{\tiny{[P1$\rangle$GM]}}
                \texttt{Yes} \\
\\ 
\texttt{**Reasoning:**} \\
\\ 
\texttt{To answer this, we must distinguish between the **observed probability** (what we see in the data) and the **counterfactual probability** (what would have happened under a different condition).} \\
\\ 
\texttt{1.  **Analyze the Observed Data:**} \\
\texttt{    *   Let \$H\$ be the event the husband sets the alarm (\$H=1\$) or does not (\$H=0\$).} \\
\texttt{    *   Let \$A\$ be the event the alarm rings (\$A=1\$) or does not (\$A=0\$).} \\
\texttt{    *   The problem states:} \\
\texttt{        *   \$P(A=1 | H=0) = 0.51\$} \\
\texttt{        *   \$P(A=1 | H=1) = 0.49\$} \\
\texttt{    *   Observationally, the alarm is *less* likely to ring when the husband sets it. This suggests a negative association between the husband setting the alarm and the alarm ringing.} \\
\\ 
\texttt{2.  **Analyze the Causal Structure:**} \\
\texttt{    *   The prompt specifies: "Husband has a direct effect on wife and alarm clock. Wife has a direct effect on alarm clock."} \\
\texttt{    *   This implies a causal graph where \$Husband \textbackslash{}rightarrow Wife\$ and \$Husband \textbackslash{}rightarrow Alarm\$, and \$Wife \textbackslash{}rightarrow Alarm\$.} \\
\texttt{    *   Crucially, the husband's action (\$H\$) influences the wife's state, which in turn influences the alarm. This creates a "fork" or a specific confounding structure where the husband's action is the common cause of both the wife's behavior and the direct alarm setting.} \\
\\ 
\texttt{3.  **Evaluate the Counterfactual:**} \\
\texttt{    *   The question asks: "For husbands that set the alarm (\$H=1\$), would it be more likely to see ringing alarm if the husband had not set the alarm (\$H=0\$)?"} \\
\texttt{    *   This is asking for the counterfactual probability: \$P(A=1 | do(H=0), H=1)\$? Or more simply, comparing the outcome of the specific unit (a husband who set the alarm) under the intervention \$H=0\$ versus the observed outcome under \$H=1\$.} \\
\texttt{    *   However, in a deterministic or probabilistic causal model defined *only* by the conditions given, we look at the **Average Treatment Effect (ATE)** or the specific causal effect of \$H\$ on \$A\$.} \\
\texttt{    *   The prompt states the world is "self{-}contained" with "only the following conditions." It gives us the conditional probabilities \$P(A|H)\$. In the absence of information about unobserved confounders (which the prompt explicitly excludes by saying "without any unmentioned factors"), the observed conditional probability \$P(A|H)\$ is the best estimate of the causal effect \$P(A|do(H))\$} \\
            }
        }
    }
    & & \\ \\

    \theutterance \stepcounter{utterance}  
    & & & \multicolumn{2}{p{0.3\linewidth}}{
        \cellcolor[rgb]{0.9,0.9,0.9}{
            \makecell[{{p{\linewidth}}}]{
                \texttt{\tiny{[GM$|$GM]}}
                \texttt{Yes} \\
            }
        }
    }
    & & \\ \\

    \theutterance \stepcounter{utterance}  
    & & & \multicolumn{2}{p{0.3\linewidth}}{
        \cellcolor[rgb]{0.9,0.9,0.9}{
            \makecell[{{p{\linewidth}}}]{
                \texttt{\tiny{[GM$|$GM]}}
                \texttt{yes} \\
            }
        }
    }
    & & \\ \\

    \theutterance \stepcounter{utterance}  
    & & & \multicolumn{2}{p{0.3\linewidth}}{
        \cellcolor[rgb]{0.9,0.9,0.9}{
            \makecell[{{p{\linewidth}}}]{
                \texttt{\tiny{[GM$|$GM]}}
                \texttt{game\_result = WIN} \\
            }
        }
    }
    & & \\ \\

\end{supertabular}
}

\end{document}
